\documentclass[10pt]{article}
\usepackage{arydshln}
\usepackage{graphicx}
\usepackage{tcolorbox}
\usepackage[a4paper, margin=1in]{geometry}
\usepackage{amsmath, amsfonts, amssymb}
\usepackage[colorlinks=true,linkcolor=black,citecolor=black,urlcolor=black]{hyperref}
\usepackage{natbib}
\usepackage{url}

% Citing how to
% \nocite{*}
% \bibliographystyle{plain}
% \bibliography{refs}
% \cite{ref of cite}

% ---- Macros ----
\newcommand{\myhrule}[1]{%
  \vspace{#1}%
  \noindent\hrule
  \vspace{#1}%
}
\makeatletter
\newcommand{\bern}{%
  \begin{enumerate}%
  \renewcommand{\labelenumi}{\roman{enumi}.}%
}
\newcommand{\eern}{%
  \end{enumerate}%
}
\makeatother
\newcommand{\FlaThreeByThreeTL}[9]{%
\left(
\begin{array}{cc|c}
#1 & #2 & #3\\
#4 & #5 & #6\\
\hline
#7 & #8 & #9
\end{array}
\right)%
}
\tcbuselibrary{breakable}
\newcommand{\boxedtext}[1]{%
\begin{tcolorbox}[
    colback=white,
    colframe=black,
    sharp corners,
    boxrule=0.5pt,
    breakable
]%
#1
\end{tcolorbox}%
}
\newcommand{\bt}[1]{%
\begin{tcolorbox}[
    colback=white,
    colframe=black,
    sharp corners,
    boxrule=0.5pt,
    breakable
]%
#1
\end{tcolorbox}%
}
\newcommand{\FlaTwoByOne}[2]{%
\left(
\begin{array}{c}
#1\\
\hline
#2
\end{array}
\right)%
}
\newcommand{\FlaThreeByOneB}[3]{%
\left(
\begin{array}{c}
#1\\
\hline
#2\\
#3
\end{array}
\right)%
}
\newcommand{\FlaTwoByTwoSingleLine}[4]{%
\left(
\begin{array}{cc}
#1 & #2\\
#3 & #4
\end{array}
\right)%
}
\newcommand{\complexpm}{\boxed{\pm}}
\newcommand{\FlaOneByTwoSingleLine}[2]{%
\left(
\begin{array}{cc}
    #1 & #2
\end{array}
\right)%
}
\newcommand{\FlaTwoByOneSingleLine}[2]{%
\left(
\begin{array}{cc}
#1\\
#2
\end{array}
\right)%
}
\newcommand{\FlaThreeByThreeBR}[9]{%
\left(
\begin{array}{c|cc}
#1 & #2 & #3 \\ \hline
#4 & #5 & #6 \\
#7 & #8 & #9
\end{array}
\right)%
}
\newcommand{\FlaTwoByTwo}[4]{%
\left(
\begin{array}{c|c}
#1 & #2 \\
\hline
#3 & #4
\end{array}
\right)
}
\newcommand{\FlaOneByThreeR}[3]{%
\left(
\begin{array}{c|cc}
#1 & #2 & #3
\end{array}
\right)%
}
\newcommand{\FlaOneByTwo}[2]{%
\left(
\begin{array}{c|c}
#1 & #2
\end{array}
\right)%
}
\newcommand{\Span}{\operatorname{Span}}
\newcommand{\pbreak}{\vspace{30pt}}
\newcommand{\ud}[1]{\underline{#1}}
\newcommand{\ital}[1]{\textit{#1}}
\newcommand{\fb}[1]{\fbox{#1}}
\newcommand{\vocab}[1]{\fbox{\textbf{\textit{#1}}}}
\newcommand{\eps}{\varepsilon}
\newcommand{\mb}{\mathbf}
\newcommand{\bs}{\boldsymbol}
\newcommand{\mdef}{\overset{\mathrm{def}}{=}}
\newcommand{\h}[1]{\textbf{\fbox{#1}}}
\newcommand{\real}{\mathbb{R}}
\newcommand{\R}{\mathbb{R}}
\newcommand{\C}{\mathbb{C}}
\newcommand{\Null}{\mathcal{N}}
\newcommand{\amp}{&}
\newcommand{\tpose}{^{\top}}
\newcommand{\ident}[1][]{\mathbf{I}_{#1}}
\newcommand{\defeq}{\overset{\cdot}{=}}
\newcommand{\sbsi}[4]{%
	\begin{center}
   		\begin{minipage}{0.5\textwidth}
      			\centering
      			\includegraphics[width=#2\textwidth]{#1}
    		\end{minipage}%
    		\begin{minipage}{0.5\textwidth}
      			\centering
      			\includegraphics[width=#4\textwidth]{#3}
		\end{minipage}
 	\end{center}
}
\newcommand{\ci}[2]{%
  \begin{center}
	\includegraphics[width=#2\textwidth]{#1}
  \end{center}
}
\newcommand{\cii}[4]{%
  \begin{center}
    \includegraphics[width=#2\textwidth]{#1}\hfill
    \includegraphics[width=#4\textwidth]{#3}
  \end{center}
}
\newcommand{\newbp}{%
  \newpage
  \mbox{}
  \newpage
}
\newcommand{\formbox}[1]{%
  \begin{center}
    \fbox{\parbox{0.95\linewidth}{\[
      #1
    \]}}
  \end{center}
}
\newcommand{\euler}{\mathrm{e}}
\newcommand{\expect}{\mathbb{E}}
\newcommand{\bi}{\begin{itemize}}
\newcommand{\ei}{\end{itemize}}
\newcommand{\be}{\begin{enumerate}}
\newcommand{\ee}{\end{enumerate}}
\newcommand{\al}[1]{\ensuremath{\begin{aligned}#1\end{aligned}}}
\DeclareMathOperator*{\argmax}{arg\,max}
\DeclareMathOperator*{\argmin}{arg\,min}
\newcommand{\rightarrowname}[1]{\ensuremath{\xrightarrow{\text{#1}}}}
\newcommand{\leftarrowname}[1]{\ensuremath{\xleftarrow{\text{#1}}}}
\newcommand{\lt}{<}
\newcommand{\gt}{>}
\setlength{\parindent}{0pt}
\newcommand{\tab}{\hspace*{2em}}
\newcommand{\pause}{}
% ------------------------------
\DeclareMathSizes{10}{13}{8}{8}

% optionally:
% \setcounter{secnumdepth}{0} % removes section numbers
\renewcommand{\arraystretch}{1.5}
% \DeclareMathSizes{<text size>}{<math size>}{<script size>}{<scriptscript size>}

\title{PostgreSQL notes}
\author{}
\date{}

\begin{document}

\maketitle

\tableofcontents

\newpage
\section{Commands}
\newpage
\section{PostgreSQL Commands (Comprehensive Cheat Sheet)}

% ================= SELECT =================
\newpage
\subsection{SELECT}
\begin{itemize}
    \item Purpose: Retrieve data from one or more tables.
\end{itemize}

\begin{verbatim}
SELECT * FROM employees;

SELECT name, salary
FROM employees
WHERE salary > 50000;
\end{verbatim}

% ================= WHERE =================
\newpage
\subsection{WHERE}
\begin{itemize}
    \item Purpose: Filter rows based on conditions.
\end{itemize}

\begin{verbatim}
SELECT *
FROM employees
WHERE department = 'Engineering'
  AND salary > 70000;
\end{verbatim}

% ================= ORDER BY =================
\newpage
\subsection{ORDER BY}
\begin{itemize}
    \item Purpose: Sort results ascending or descending.
\end{itemize}

\begin{verbatim}
SELECT *
FROM employees
ORDER BY salary DESC;
\end{verbatim}

% ================= DISTINCT =================
\newpage
\subsection{DISTINCT}
\begin{itemize}
    \item Purpose: Remove duplicates.
\end{itemize}

\begin{verbatim}
SELECT DISTINCT department
FROM employees;
\end{verbatim}

% ================= LIMIT / OFFSET =================
\newpage
\subsection{LIMIT / OFFSET}
\begin{itemize}
    \item Purpose: Restrict number of rows and paginate results.
\end{itemize}

\begin{verbatim}
SELECT *
FROM employees
LIMIT 10 OFFSET 20;
\end{verbatim}

% ================= GROUP BY =================
\newpage
\subsection{GROUP BY}
\begin{itemize}
    \item Purpose: Aggregate rows into groups.
\end{itemize}

\begin{verbatim}
SELECT department, COUNT(*)
FROM employees
GROUP BY department;
\end{verbatim}

% ================= HAVING =================
\newpage
\subsection{HAVING}
\begin{itemize}
    \item Purpose: Filter aggregated groups.
\end{itemize}

\begin{verbatim}
SELECT department, COUNT(*)
FROM employees
GROUP BY department
HAVING COUNT(*) > 5;
\end{verbatim}

% ================= JOINS =================
\newpage
\subsection{INNER JOIN}
\begin{itemize}
    \item Purpose: Only matching rows between tables.
\end{itemize}

\begin{verbatim}
SELECT e.name, d.name
FROM employees e
INNER JOIN departments d
ON e.department_id = d.id;
\end{verbatim}

\newpage
\subsection{LEFT JOIN}
\begin{itemize}
    \item Purpose: All left table rows + matches from right.
\end{itemize}

\begin{verbatim}
SELECT e.name, d.name
FROM employees e
LEFT JOIN departments d
ON e.department_id = d.id;
\end{verbatim}

\newpage
\subsection{RIGHT JOIN}
\begin{itemize}
    \item Purpose: All right table rows + matches from left.
\end{itemize}

\begin{verbatim}
SELECT e.name, d.name
FROM employees e
RIGHT JOIN departments d
ON e.department_id = d.id;
\end{verbatim}

\newpage
\subsection{FULL OUTER JOIN}
\begin{itemize}
    \item Purpose: All rows from both tables.
\end{itemize}

\begin{verbatim}
SELECT e.name, d.name
FROM employees e
FULL OUTER JOIN departments d
ON e.department_id = d.id;
\end{verbatim}

% ================= INSERT =================
\newpage
\subsection{INSERT}
\begin{itemize}
    \item Purpose: Insert new rows.
\end{itemize}

\begin{verbatim}
INSERT INTO employees (name, salary)
VALUES ('Alice', 90000);
\end{verbatim}

% ================= UPDATE =================
\newpage
\subsection{UPDATE}
\begin{itemize}
    \item Purpose: Modify existing rows.
\end{itemize}

\begin{verbatim}
UPDATE employees
SET salary = salary * 1.1
WHERE department = 'Engineering';
\end{verbatim}

% ================= DELETE =================
\newpage
\subsection{DELETE}
\begin{itemize}
    \item Purpose: Remove rows.
\end{itemize}

\begin{verbatim}
DELETE FROM employees
WHERE id = 10;
\end{verbatim}

% ================= CASE =================
\newpage
\subsection{CASE WHEN}
\begin{itemize}
    \item Purpose: Conditional logic in queries.
\end{itemize}

\begin{verbatim}
SELECT name,
       CASE
           WHEN salary > 100000 THEN 'High'
           WHEN salary > 50000 THEN 'Medium'
           ELSE 'Low'
       END AS salary_level
FROM employees;
\end{verbatim}

% ================= IN =================
\newpage
\subsection{IN}
\begin{itemize}
    \item Purpose: Match against multiple values.
\end{itemize}

\begin{verbatim}
SELECT *
FROM employees
WHERE department IN ('Engineering', 'HR');
\end{verbatim}

% ================= BETWEEN =================
\newpage
\subsection{BETWEEN}
\begin{itemize}
    \item Purpose: Range filtering.
\end{itemize}

\begin{verbatim}
SELECT *
FROM employees
WHERE salary BETWEEN 50000 AND 100000;
\end{verbatim}

% ================= LIKE =================
\newpage
\subsection{LIKE}
\begin{itemize}
    \item Purpose: Pattern matching.
\end{itemize}

\begin{verbatim}
SELECT *
FROM employees
WHERE name LIKE 'A%';
\end{verbatim}

% ================= NULL CHECK =================
\newpage
\subsection{NULL CHECK}
\begin{itemize}
    \item Purpose: Check missing values.
\end{itemize}

\begin{verbatim}
SELECT *
FROM employees
WHERE manager_id IS NULL;
\end{verbatim}

% ================= AGGREGATES =================
\newpage
\subsection{Aggregate Functions}
\begin{itemize}
    \item Purpose: Compute summary statistics.
\end{itemize}

\begin{verbatim}
SELECT
    COUNT(*),
    SUM(salary),
    AVG(salary),
    MAX(salary),
    MIN(salary)
FROM employees;
\end{verbatim}

% ================= WINDOW FUNCTIONS =================
\newpage
\subsection{Window Functions}
\begin{itemize}
    \item Purpose: Compute values across partitions without collapsing rows.
\end{itemize}

\begin{verbatim}
SELECT name,
       salary,
       RANK() OVER (PARTITION BY department ORDER BY salary DESC)
FROM employees;
\end{verbatim}

% ================= CTE =================
\newpage
\subsection{WITH (CTE)}
\begin{itemize}
    \item Purpose: Create temporary named result sets.
\end{itemize}

\begin{verbatim}
WITH high_salary AS (
    SELECT *
    FROM employees
    WHERE salary > 100000
)
SELECT * FROM high_salary;
\end{verbatim}

% ================= SUBQUERY =================
\newpage
\subsection{Subqueries}
\begin{itemize}
    \item Purpose: Nest queries inside another query.
\end{itemize}

\begin{verbatim}
SELECT name
FROM employees
WHERE salary > (
    SELECT AVG(salary)
    FROM employees
);
\end{verbatim}

% ================= UNION =================
\newpage
\subsection{UNION / UNION ALL}
\begin{itemize}
    \item Purpose: Combine results of multiple queries.
\end{itemize}

\begin{verbatim}
SELECT name FROM employees
UNION
SELECT name FROM contractors;
\end{verbatim}

% ================= EXISTS =================
\newpage
\subsection{EXISTS}
\begin{itemize}
    \item Purpose: Check existence of rows in subquery.
\end{itemize}

\begin{verbatim}
SELECT name
FROM employees e
WHERE EXISTS (
    SELECT 1
    FROM departments d
    WHERE d.id = e.department_id
);
\end{verbatim}

% ================= INDEX =================
\newpage
\subsection{INDEX}
\begin{itemize}
    \item Purpose: Speed up queries.
\end{itemize}

\begin{verbatim}
CREATE INDEX idx_salary
ON employees(salary);
\end{verbatim}

% ================= TRANSACTIONS =================
\newpage
\subsection{Transactions}
\begin{itemize}
    \item Purpose: Ensure atomic operations.
\end{itemize}

\begin{verbatim}
BEGIN;

UPDATE employees SET salary = salary * 1.1;

COMMIT;
\end{verbatim}

\newpage 
\subsection{ALTER TABLE}
\subsection{DROP TABLE}
\subsection{TRUNCATE}
\subsection{CREATE TABLE}
\subsection{CREATE DATABASE}
\subsection{CREATE VIEW}
\subsection{DROP VIEW}
\subsection{INSERT INTO ... ON CONFLICT}
\subsection{UPSERT (ON CONFLICT DO UPDATE)}
\subsection{MERGE}
\subsection{RETURNING}
\subsection{COALESCE}
\subsection{NULLIF}
\subsection{CAST}
\subsection{:: type casting}
\subsection{STRING FUNCTIONS (CONCAT, SUBSTRING, LENGTH)}
\subsection{DATE FUNCTIONS (NOW, CURRENT\_DATE, AGE)}
\subsection{INTERVAL}
\subsection{ARRAYS}
\subsection{UNNEST}
\subsection{JSON / JSONB OPERATIONS}
\subsection{CTE RECURSIVE}
\subsection{WINDOW FUNCTIONS (LEAD)}
\subsection{WINDOW FUNCTIONS (LAG)}
\subsection{NTILE}
\subsection{ROW\_NUMBER}
\subsection{RANK}
\subsection{DENSE\_RANK}
\subsection{PARTITION BY}
\subsection{OVER CLAUSE}
\subsection{EXPLAIN}
\subsection{EXPLAIN ANALYZE}
\subsection{VACUUM}
\subsection{ANALYZE}
\subsection{REINDEX}
\subsection{GRANT}
\subsection{REVOKE}

\newbp 
\section{General}

\subsection{General Structure of a SQL Query}
\begin{itemize}
\item A SQL query follows a logical execution order, which is different from how it is written.
\end{itemize}

\begin{verbatim}
SELECT columns
FROM table
JOIN other_table
ON join_condition
WHERE row_filter
GROUP BY grouping_columns
HAVING group_filter
ORDER BY sorting_columns
LIMIT number_of_rows;
\end{verbatim}

\begin{itemize}
\item Logical execution order:
    \begin{itemize}
    \item FROM (choose tables)
    \item JOIN (combine tables)
    \item WHERE (filter rows)
    \item GROUP BY (group rows)
    \item HAVING (filter groups)
    \item SELECT (choose columns / expressions)
    \item ORDER BY (sort results)
    \item LIMIT (restrict output)
    \end{itemize}
\item Important: The written order is NOT the execution order.
\end{itemize}

\newpage
\subsection{Join Differences}
The matches will fill in a NULL value for entries that don't have a match\\ 
(in the case of like all left or all right).
\begin{itemize}
\item INNER JOIN -- only matches
\item LEFT JOIN -- all left + matches
\item RIGHT JOIN -- all right + matches
\item FULL JOIN -- all rows from both sides
\end{itemize}

\subsection{NULL Behavior}
NULL means "unknown value" and behaves differently from normal values.
\begin{itemize}
\item Use IS NULL / IS NOT NULL (NOT =)
\item NULL in comparisons results in UNKNOWN (not true/false)
\item Aggregates ignore NULL values (except COUNT(*))
\end{itemize}

\subsection{WHERE vs HAVING}
\begin{itemize}
\item WHERE filters rows BEFORE grouping
\item HAVING filters AFTER GROUP BY
\end{itemize}

\subsection{GROUP BY Rule}
\begin{itemize}
\item Every selected column must be either:
    \begin{itemize}
    \item in GROUP BY
    \item or an aggregate function (COUNT, MAX, etc.)
    \end{itemize}
\end{itemize}

\subsection{ORDER BY vs GROUP BY}
\begin{itemize}
\item GROUP BY = aggregates rows
\item ORDER BY = sorts final result
\end{itemize}

\subsection{COUNT Differences}
\begin{itemize}
\item COUNT(*) counts all rows
\item COUNT(column) ignores NULL values
\end{itemize}

\subsection{DISTINCT}
\begin{itemize}
\item Removes duplicate rows from result
\item Works on entire selected row combination
\end{itemize}

\subsection{JOIN Condition}
\begin{itemize}
\item ON clause defines how tables match
\item Missing or wrong ON can cause Cartesian product (explosion of rows)
\end{itemize}

\subsection{Cartesian Product}
\begin{itemize}
\item Occurs when JOIN has no condition
\item Every row combines with every other row
\item Avoid unless intentionally using CROSS JOIN
\end{itemize}

\subsection{Aliases}
\begin{itemize}
\item Temporary names for tables or columns
\item Improves readability and reduces repetition
\end{itemize}

\end{document}


